% GENERATED FILE. Edit canonical lessons, then run npm run generate:books.
% canonical-source-hash: fnv1a64:74477d86f513d756
% canonical-lessons: SA-C475-svidyati, SA-C475-jrmbhate, SA-C475-ksauti, SA-C475-kanduyati, SA-C475-hikka

\chapter{Sweat, Yawns and Sneezes}
\label{ch:sa-sweat-yawns-and-sneezes}
\hlchaptermodality{475}

\begin{chapteropening}
\textbf{By the end of this chapter:} \emph{I can say what the body does by itself.}

Complete the last lesson of chapter 475: I can say what the body does by itself.
\end{chapteropening}

\section[\texorpdfstring{\sk{स्विद्यति} (svidyati)}{स्विद्यति (svidyati)}]{\sk{स्विद्यति} (svidyati) — sweats}

\label{lesson:SA-C475-svidyati}



\begin{quote}
Before the new one: say the Sanskrit for flesh, then the Sanskrit for throbs.
\end{quote}

\subsection*{You'll want to know: \sk{स्विद्यति}}
\textbf{\sk{स्विद्यति}} — \emph{svidyati} — \textquotedblleft{}sweats\textquotedblright{}.

\begin{quote}
\sk{ग्रीष्मे} \sk{कृषकः} \sk{बहु} \sk{स्विद्यति।}
\end{quote}

\begin{quote}
\emph{grīṣme kṛṣakaḥ bahu svidyati.}
\end{quote}

\begin{quote}
In summer the farmer sweats a lot.
\end{quote}

Past (he …): \sk{सः} \sk{बहु} \sk{स्विन्नवान्}, \emph{saḥ bahu svinnavān}.

\begin{cousinweb}
\textbf{\sk{स्विद्यति}} (\emph{svidyati}) comes from the root √svid, \textquotedblleft{}to sweat, to perspire\textquotedblright{}; MW's ppp \emph{svinna} is \textquotedblleft{}sweating\textquotedblright{} (MW).
\end{cousinweb}

\begin{grammarlens}[title={What you've built}]
One more word for describing what you lived and felt.
\end{grammarlens}

\subsection*{Your turn}
\begin{itemize}
\raggedright
  \item \emph{Say it:} \emph{svidyati}
  \item \emph{Say it:} \emph{svidyati}, once more
  \item \emph{Say it:} \emph{grīṣme kṛṣakaḥ bahu svidyati.} — In summer the farmer sweats a lot.
  \item \emph{Recall:} say the Sanskrit for flesh, then the Sanskrit for throbs, then say \emph{svidyati} again
\end{itemize}

\begin{tcolorbox}[breakable,colback=teal!4,colframe=teal!35!black,title={Before you move on}]
What does \sk{स्विद्यति} mean? (\textquotedblleft{}Sweats\textquotedblright{}.) Say it once more.
\end{tcolorbox}
\section[\texorpdfstring{\sk{जृम्भते} (jṛmbhate)}{जृम्भते (jṛmbhate)}]{\sk{जृम्भते} (jṛmbhate) — yawns}

\label{lesson:SA-C475-jrmbhate}



\begin{quote}
Before the new one: say the Sanskrit for throbs, then the Sanskrit for sweats.
\end{quote}

\subsection*{You'll want to know: \sk{जृम्भते}}
\textbf{\sk{जृम्भते}} — \emph{jṛmbhate} — \textquotedblleft{}yawns\textquotedblright{}.

\begin{quote}
\sk{रात्रौ} \sk{श्रान्तः} \sk{बालः} \sk{जृम्भते।}
\end{quote}

\begin{quote}
\emph{rātrau śrāntaḥ bālaḥ jṛmbhate.}
\end{quote}

\begin{quote}
At night the tired child yawns.
\end{quote}

Past (he …): \sk{सः} \sk{जृम्भितवान्}, \emph{saḥ jṛmbhitavān}.

\begin{cousinweb}
\textbf{\sk{जृम्भते}} (\emph{jṛmbhate}) comes from the root √jṛmbh, \textquotedblleft{}to open the mouth, to yawn\textquotedblright{} (MW).
\end{cousinweb}

\begin{grammarlens}[title={What you've built}]
One more word for describing what you lived and felt.
\end{grammarlens}

\subsection*{Your turn}
\begin{itemize}
\raggedright
  \item \emph{Say it:} \emph{jṛmbhate}
  \item \emph{Say it:} \emph{jṛmbhate}, once more
  \item \emph{Say it:} \emph{rātrau śrāntaḥ bālaḥ jṛmbhate.} — At night the tired child yawns.
  \item \emph{Recall:} say the Sanskrit for throbs, then the Sanskrit for sweats, then say \emph{jṛmbhate} again
\end{itemize}

\begin{tcolorbox}[breakable,colback=teal!4,colframe=teal!35!black,title={Before you move on}]
What does \sk{जृम्भते} mean? (\textquotedblleft{}Yawns\textquotedblright{}.) Say it once more.
\end{tcolorbox}
\section[\texorpdfstring{\sk{क्षौति} (kṣauti)}{क्षौति (kṣauti)}]{\sk{क्षौति} (kṣauti) — sneezes}

\label{lesson:SA-C475-ksauti}



\begin{quote}
Before the new one: say the Sanskrit for sweats, then the Sanskrit for yawns.
\end{quote}

\subsection*{You'll want to know: \sk{क्षौति}}
\textbf{\sk{क्षौति}} — \emph{kṣauti} — \textquotedblleft{}sneezes\textquotedblright{}.

\begin{quote}
\sk{धूल्या} \sk{सः} \sk{क्षौति।}
\end{quote}

\begin{quote}
\emph{dhūlyā saḥ kṣauti.}
\end{quote}

\begin{quote}
The dust makes him sneeze.
\end{quote}

Past (he …): \sk{सः} \sk{क्षुतवान्}, \emph{saḥ kṣutavān}.

\begin{cousinweb}
\textbf{\sk{क्षौति}} (\emph{kṣauti}) comes from the root √kṣu, \textquotedblleft{}to sneeze\textquotedblright{}, a second-class verb whose vowel strengthens to \emph{au} (MW).
\end{cousinweb}

\begin{grammarlens}[title={What you've built}]
One more word for describing what you lived and felt.
\end{grammarlens}

\subsection*{Your turn}
\begin{itemize}
\raggedright
  \item \emph{Say it:} \emph{kṣauti}
  \item \emph{Say it:} \emph{kṣauti}, once more
  \item \emph{Say it:} \emph{dhūlyā saḥ kṣauti.} — The dust makes him sneeze.
  \item \emph{Recall:} say the Sanskrit for sweats, then the Sanskrit for yawns, then say \emph{kṣauti} again
\end{itemize}

\begin{tcolorbox}[breakable,colback=teal!4,colframe=teal!35!black,title={Before you move on}]
What does \sk{क्षौति} mean? (\textquotedblleft{}Sneezes\textquotedblright{}.) Say it once more.
\end{tcolorbox}
\section[\texorpdfstring{\sk{कण्डूयति} (kaṇḍūyati)}{कण्डूयति (kaṇḍūyati)}]{\sk{कण्डूयति} (kaṇḍūyati) — scratches}

\label{lesson:SA-C475-kanduyati}



\begin{quote}
Before the new one: say the Sanskrit for yawns, then the Sanskrit for sneezes.
\end{quote}

\subsection*{You'll want to know: \sk{कण्डूयति}}
\textbf{\sk{कण्डूयति}} — \emph{kaṇḍūyati} — \textquotedblleft{}scratches\textquotedblright{}.

\begin{quote}
\sk{शुनकः} \sk{कर्णं} \sk{कण्डूयति।}
\end{quote}

\begin{quote}
\emph{śunakaḥ karṇaṃ kaṇḍūyati.}
\end{quote}

\begin{quote}
The dog scratches its ear.
\end{quote}

Past (it …): \sk{शुनकः} \sk{कर्णं} \sk{कण्डूयितवान्}, \emph{śunakaḥ karṇaṃ kaṇḍūyitavān}.

\begin{cousinweb}
\textbf{\sk{कण्डूयति}} (\emph{kaṇḍūyati}) is built on \emph{kaṇḍū}, \textquotedblleft{}itching, the itch\textquotedblright{}; MW gives \textquotedblleft{}to scratch, rub\textquotedblright{} and \textquotedblleft{}to itch\textquotedblright{} (MW).
\end{cousinweb}

\begin{grammarlens}[title={What you've built}]
One more word for describing what you lived and felt.
\end{grammarlens}

\subsection*{Your turn}
\begin{itemize}
\raggedright
  \item \emph{Say it:} \emph{kaṇḍūyati}
  \item \emph{Say it:} \emph{kaṇḍūyati}, once more
  \item \emph{Say it:} \emph{śunakaḥ karṇaṃ kaṇḍūyati.} — The dog scratches its ear.
  \item \emph{Recall:} say the Sanskrit for yawns, then the Sanskrit for sneezes, then say \emph{kaṇḍūyati} again
\end{itemize}

\begin{tcolorbox}[breakable,colback=teal!4,colframe=teal!35!black,title={Before you move on}]
What does \sk{कण्डूयति} mean? (\textquotedblleft{}Scratches\textquotedblright{}.) Say it once more.
\end{tcolorbox}
\section[\texorpdfstring{\sk{हिक्का} (hikkā)}{हिक्का (hikkā)}]{\sk{हिक्का} (hikkā) — a hiccup}

\label{lesson:SA-C475-hikka}



\begin{quote}
Before the new one: say the Sanskrit for sneezes, then the Sanskrit for scratches.
\end{quote}

\subsection*{You'll want to know: \sk{हिक्का}}
\textbf{\sk{हिक्का}} — \emph{hikkā} — \textquotedblleft{}a hiccup\textquotedblright{}.

\begin{quote}
\sk{सः} \sk{शीघ्रं} \sk{खादति}, \sk{अतः} \sk{तस्य} \sk{हिक्का} \sk{भवति।}
\end{quote}

\begin{quote}
\emph{saḥ śīghraṃ khādati, ataḥ tasya hikkā bhavati.}
\end{quote}

\begin{quote}
He eats too fast, so he gets hiccups.
\end{quote}

A feminine noun.

\begin{cousinweb}
\textbf{\sk{हिक्का}} (\emph{hikkā}) MW gives \textquotedblleft{}hiccup, sob, a spasmodic sound in the throat\textquotedblright{} and compares the form \emph{hekkā}; MW gives no derivation, so its origin is uncertain (MW).
\end{cousinweb}

\begin{grammarlens}[title={What you've built}]
One more word for describing what you lived and felt.
\end{grammarlens}

\subsection*{Your turn}
\begin{itemize}
\raggedright
  \item \emph{Say it:} \emph{hikkā}
  \item \emph{Say it:} \emph{hikkā}, once more
  \item \emph{Say it:} \emph{saḥ śīghraṃ khādati, ataḥ tasya hikkā bhavati.} — He eats too fast, so he gets hiccups.
  \item \emph{Recall:} say the Sanskrit for sneezes, then the Sanskrit for scratches, then say \emph{hikkā} again
\end{itemize}

\begin{tcolorbox}[breakable,colback=teal!4,colframe=teal!35!black,title={Before you move on}]
What does \sk{हिक्का} mean? (\textquotedblleft{}A hiccup\textquotedblright{}.) Say it once more.
\end{tcolorbox}
